\chapter{Synthesis: The Three Faces of One Group}
\label{ch:synthesis}

We can now assemble the whole argument without skipping any major arrow.

\begin{center}
\resizebox{0.86\linewidth}{!}{\classicTrinityDiagram}
\end{center}

\begin{enumerate}
\item A lattice \(\Lambda\subset\C\) gives a compact genus-one complex Lie group
  \(T=\C/\Lambda\).
\item Elliptic functions are meromorphic functions on \(T\).
\item The Weierstrass function \(\wp\) has one double pole at \(O\), and
  \(\wp'\) has one triple pole.
\item The differential equation
  \(\wp'^2=4\wp^3-g_2\wp-g_3\) produces a cubic \(E\).
\item Riemann--Roch explains why \(1,\wp,\wp'\) are exactly the functions needed
  for the projective map \(T\to\CP^2\).
\item The map \(\Phi=[\wp:\wp':1]\) identifies \(T\) with the smooth cubic
  \(E(\C)\).
\item Divisors and Abel's theorem identify \(\Pic^0(E)\) with the same group.
\item A line meeting \(E\) in \(P,Q,R\) cuts a principal divisor
  \([P]+[Q]+[R]-3[O]\).
\item Abel's theorem gives \(P+Q+R=O\), so the group sum of \(P\) and \(Q\) is
  the reflection of \(R\).
\item Because the addition formulas are rational, the same cubic opens into
  arithmetic over \(\Q\) and finite fields.
\end{enumerate}

\begin{theorem}[Three-face theorem]
\label{thm:three-face}
Over \(\C\), the analytic torus \(\C/\Lambda\), the Jacobian \(\Pic^0(E)\), and
the smooth cubic \(E(\C)\) are isomorphic as complex Lie groups.  Equivalently,
the identifying maps are biholomorphic group isomorphisms.  Under these
identifications, the chord--tangent rule on \(E\) is the group law on
\(\C/\Lambda\) transported by \(\Phi=[\wp:\wp':1]\).
\end{theorem}

\begin{proof}
The analytic-to-algebraic identification is \cref{thm:wp-embedding-new}.  The
geometric identification with the Jacobian is \cref{cor:pic-torus} transported
through \(\Phi\).  The equality between transported addition and chord--tangent
addition is \cref{thm:transport-new}.
\end{proof}

\section*{Exercises}
\addcontentsline{toc}{section}{Exercises}

\exerciseintent{Recover the forced reflection step from the line-divisor Abel equation.}
\begin{exercise}
Compute the line-divisor relation for collinear \(P,Q,R\):
\[
  [P]+[Q]+[R]-3[O]\in\Prin(E).
\]
Apply Abel's theorem to get \(P+Q+R=O\), and solve this equation for the
chord--tangent sum \(P\boxplus Q\).
\end{exercise}

\exerciseintent{Organize the three languages and the maps between them.}
\begin{exercise}
Make a three-column table with headings analytic, divisor-theoretic, and
algebraic.  Place \(\C/\Lambda\), \(\Pic^0(E)\), \(E(\C)\), \(\wp\), divisors,
and lines in the appropriate columns.  Add a second row containing the explicit
maps
\[
  \Phi=[\wp:\wp':1],\qquad P\mapsto[P]-[O],
\]
and the Abel--Jacobi map.
\end{exercise}

\exerciseintent{Map each chapter block to the calculation type it contributes.}
\begin{exercise}
Make a chapter-to-method table.  For each block of chapters, record one
calculation type: residue or integral calculation, lattice linear-algebra
calculation, pole-order dimension count, divisor coefficient calculation, line
intersection calculation, cohomology quotient calculation, or finite-field point
count.
\end{exercise}

\exerciseintent{Audit the two main maps as complex Lie group isomorphisms and mark theorem inputs.}
\begin{exercise}
Make a map-check table for
\[
  \C/\Lambda\cong\Pic^0(E)\cong E(\C)
\]
as complex Lie groups.  Use rows \(\Phi\) and \(P\mapsto[P]-[O]\), and columns:
formula, source, target, group-law compatibility, holomorphicity input, and
inverse or bijectivity input.  In each entry, mark whether it is a direct
calculation from the exercises or an accepted named theorem.
\end{exercise}

\exerciseintent{Compute final Jacobian addition from the Abel sum of \([P_1]+[P_2]-[S]-[O]\).}
\begin{exercise}
Give a final dictionary for a point \(z\in\C/\Lambda\).  Include its torus
notation, its cubic notation, its divisor-class notation, and the rule for
adding two such points in each language.  Write the formulas for
\(\Phi(z)\), \([\Phi(z)]-[O]\), \(\Phi(z_1+z_2)\), and the corresponding sum in
\(\Pic^0(E)\).  Then, for \(P_i=\Phi(z_i)\) and \(S=\Phi(z_1+z_2)\), compute
the Abel sum of
\[
  [P_1]+[P_2]-[S]-[O]
\]
and use Abel's theorem to prove
\[
  ([P_1]-[O])+([P_2]-[O])=[S]-[O]
\]
in \(\Pic^0(E)\).
\end{exercise}
